\begin{helper}[Separation from an elementary prefix's missing witnesses]
Let $G$ be a finite loopless multigraph with individual edge identities,
let $\psi$ be a proper partial coloring with $k$ colors, and let $a\ne b$
be colors. Let $B\subseteq V(G)$ be elementary under $\psi$.
Suppose $u,v\in B$, $w\ne u,v$, $a$ is missing at $u$, $b$ is missing
at $v$, and at least one of $a,b$ is missing at $w$. Suppose the
$a,b$ components rooted at $u$ and $v$ are equal. Then the component
rooted at $w$ is disjoint from that common component and is a
\noderef{TwoColorPathComponent}, including the isolated-vertex case.
Moreover, swapping this component preserves the entire missing-color
set at every vertex of $B$.
\end{helper}
\begin{proof}
Write $J=\noderef{TwoColorIncidenceGraph}(G,\psi,a,b)$.
\noderef{MissingColorIncidenceEndpoint} implies that each vertex missing
one affected color has at most one neighbor in $J$; for a vertex missing
$b$, interchange $a,b$ in that result, which leaves $J$ unchanged.
If $u=v$, that vertex misses both colors. The definition of $J$ then
gives no incident adjacency at its vertex-node, so its component is a
singleton. It cannot contain $w$.

If $u\ne v$, suppose $w$ belongs to their common component.
\noderef{TwoColorComponentSpanningPath} supplies a simple path whose
support is this component: the component has a degree-at-most-one
vertex-node, namely $u$, so it is not a cycle. Any vertex-node on a
simple path with at most one neighbor must be one of its two endpoints:
an internal occurrence has distinct predecessor and successor neighbors.
Thus $u,v,w$ would be three distinct endpoints of a single path, which
is impossible. In either case $w$ is outside the common component.
By \noderef{TwoColorComponentsDisjoint}, components are either equal
or disjoint. Equality would put $w$ in that component, so they are
disjoint. The vertex-node $w$ belongs to its own component and has at
most one neighbor; these facts and $a\ne b$ are the fields of
\noderef{TwoColorPathComponent}.

Let $Q$ be the component rooted at $w$ and $\tau$ its swap. By
disjointness, $u,v$ are outside $Q$. Fix $x\in B$. If $x$ is outside
$Q$, the outside-vertex clause of
\noderef{TwoColorComponentSwapMissingBehavior} gives equality of its
missing sets. If $x$ is in $Q$, then $x\ne u,v$. Elementarity, as
defined by \noderef{ColoredVertexSetElementary}, prevents $a$ being
missing at both $x$ and $u$, and prevents $b$ being missing at both
$x$ and $v$. Thus both colors are present at $x$. Expand
\noderef{MissingColors} and \noderef{PresentColors} to choose an
incident $a$-edge and an incident $b$-edge. Each edge-node is adjacent
in $J$ to $x$, hence belongs to $Q$. The definition of
\noderef{TwoColorComponentSwap} changes the first edge to $b$ and the
second to $a$, so both colors remain present at $x$. Every other color
has the same missing predicate by the unaffected-color clause of
\noderef{TwoColorComponentSwapMissingBehavior}. Equality holds for
each color and hence for the entire missing set. This does not assert
that displayed edges at internal intersections with $B$ are fixed.
\end{proof}
